\documentclass[preprint]{revtex4-1}
\usepackage[latin1]{inputenc}
\usepackage{amsfonts}
\usepackage[T1]{fontenc}
\usepackage{xcolor}
\usepackage{color}
\usepackage{amsmath}
\usepackage{amssymb}
\usepackage{amsthm}
\usepackage{mathrsfs}
\usepackage{graphicx}
\usepackage{fancyhdr}
\usepackage{array}
\usepackage{simplewick}
\usepackage{latexsym}
\usepackage[all]{xy}
\usepackage{eufrak}
\usepackage{euscript}
\usepackage{enumerate}
%\usepackage{dsfont}
\usepackage{slashed}
\usepackage{hyperref}
\usepackage{bm}

\newcommand{\ki}[1]{\textcolor{green!60!black}{KI: #1}}

\begin{document}

\begin{center}

\Large{Revised Manuscript DS13778 Soares Rocha
}\\

\vspace{0.5cm}

{\bf \Large{Modelling stochastic fluctuations in relativistic kinetic theory}}\\

\vspace{0.5cm}

\normalsize{Gabriel Soares Rocha, Lorenzo Gavassino, Nicki Mullins}\\

\vspace{0.5cm}



\end{center}

\vspace{0.5cm}

\noindent
Dear Editor,\\

Below we respond to the remarks made by the Referee. Following their comments, new text has been added to the conclusion and section II G. The changes have been highlighted in red. We hope that after these modifications the manuscript is suitable for publication in Physical Review D. \\
 
\noindent
Best regards,\\

\noindent
The Authors\\

\newpage

\section*{Reply to the referee}

We thank the referee for carefully reading our manuscript and their constructive feedback. In the following, we reply to all the comments made by them.\\


\textit{This paper explores fluctuations in relativistic kinetic theory through thermodynamic arguments and the information current, demonstrating that the linearized Boltzmann equation ensures stability, causality, and symmetric hyperbolicity. The authors analyze significant fluctuations in small systems, such as heavy-ion collisions, and identify methods to determine correlation functions, highlighting the utility of the information current for defining a stable stochastic theory. This is a solid and important work. I only have several comments or suggestions before publication:}

\textit{1. It would be beneficial if the authors could provide a more intuitive description of the zn and zpi in Eq. (34) on page 9.
}

\textbf{R:} Fixed a generic observable $A$, the quantity $z_A$ measures the relative error that we commit when we estimate the fluctuations of $A$ using the Israel-Stewart theory instead of kinetic theory. In other words, $z_A$ is the percentage of fluctuations of $A$ that are not accounted for by relativistic hydrodynamics. We have added a clarification below Eq.~(34), in the new version.\\



\textit{2. It seems there is a typo in the first line below Eq.~(12) on page 5.
}

\textbf{R:} The text surrounding the reference below Eq.~(12) has been changed to render it more clear.\\

\textit{3. It would be great if the authors could provide some suggestions or
comments about the applications to small systems in heavy-ion collisions.}

\textbf{R:} We thank the referee for the suggestion. A paragraph has been added in the conclusion section to address this comment (the second-to-last one in the new version). Indeed, there have been some attempts to take noise effects into account in hydrodynamic simulations. However, not the codes in the state of the art of heavy-ion phenomenology. Fluctuating relativistic kinetic theory might serve as a guide to noise effects in heavy-ion phenomenology as much as regular relativistic kinetic theory has been to relativistic dissipative hydrodynamics.  

\end{document}